\section{刚体力学 II}

\subsection{定轴转动定理}

\subsubsection{角动量守恒}

设刚体关于转动轴所受力矩为 $\vect{M}$，转动惯量为 $I$，角加速度为 $\vect{\beta}$，那么有：
$$
\vect{M} = I \vect{\beta}
$$

当刚体不受合外力矩时，角动量守恒。

\subsubsection{转动动能}

刚体转动时，具有一个动能：
$$
E_k = \int \frac{1}{2} (\vect{r} \times \vect{\omega})^2 \dd{m} = \frac{1}{2} I \vect{\omega}^2
$$

\subsubsection{力矩的功}

对于一个力矩 $\vect{M}$，在其作用下刚体转动了角度 $\dd{\phi}$，其做功为：
$$
\dd{W} = \vect{M} \cdot \dd{\vect{\phi}}
$$

\subsubsection{刚体的动能合重力势能}

对于刚体，其平动动能为：
$$
E_k = \int \frac{1}{2} v^2 \dd{m} = \frac{1}{2} M v_c^2
$$
其重力势能为：
$$
E_p = \int \vect{g} \cdot \vect{r} \dd{m} = M \vect{g} \cdot \vect{r}_c
$$
